% SLIDE 8 — Pipeline de branches
\begin{frame}{Pipeline de branches (fusion)}
  \timing{1:30}
  \textbf{Branche} = un jeu de features alimentant les classifieurs.
  \textbf{\NBranches{} branches} testées :
  \vspace{0.4em}

  \centering
  \begin{tikzpicture}[
    node distance=0.35cm and 0.25cm,
    raw/.style={draw=accent2, fill=accent2!10, rounded corners, minimum width=1.6cm,
                minimum height=0.55cm, font=\scriptsize, align=center},
    br/.style={draw=accent, fill=accent!8, rounded corners, minimum width=1.5cm,
               minimum height=0.5cm, font=\scriptsize, align=center},
    arr/.style={-{Stealth[length=1.5mm]}, thick, greytx}
  ]
    \node[raw] (raw) {IRM\\prétraitée};
    \node[br, below left=0.6cm and 2.8cm of raw] (a) {A — Stats};
    \node[br, below left=0.6cm and 1.4cm of raw] (b) {B — Tex(opt)};
    \node[br, below=0.6cm of raw] (c) {C — DWT(opt)};
    \node[br, below right=0.6cm and 1.4cm of raw] (d) {D — HOG};
    \node[br, below right=0.6cm and 2.8cm of raw] (f) {Full(opt)};
    \node[raw, below=1.8cm of c, minimum width=3.2cm] (clf) {\NModels{} classifieurs};
    \foreach \x in {a,b,c,d,f} {
      \draw[arr] (raw) -- (\x);
      \draw[arr] (\x) -- (clf);
    }
    \node[font=\tiny, greytx, below=0.1cm of a] {+A+B, B+C, A+B+C};
  \end{tikzpicture}

  \vspace{0.3em}
  \small On ne devine pas la meilleure combinaison --- on construit les 8 branches et on laisse les résultats décider.
\end{frame}

\note{
On ne devine pas quelle combinaison est la meilleure --- on construit les 8 branches
et on laisse les résultats décider.
}
